\chapter{Նախագծի ընդհանուր գաղափարը}

Նպատակն է մշակել ինտեգրված ապարատա-ծրագրային հարթակ, որտեղ FPGA-ն գեներացնում կամ
ստանում է սենսորանման ազդանշաններ, կատարում է իրական ժամանակի նախնական մշակում
(preprocessing) և ստացված տվյալները փոխանցում սերվերին։ Web հարթակը պահպանում,
վիզուալիզացնում, վերլուծում և մոդելավորում է հավաքագրված տվյալները, ինչպես նաև
աշխատանքային ռեժիմում կոնֆիգուրացնում է FPGA ապարատը։

\problemstatement{Անհրաժեշտ է կառուցել \emph{Web հարթակ $\leftrightarrow$ Server
$\leftrightarrow$ Zynq PS $\leftrightarrow$ FPGA} ամբողջական համակարգ, որում FPGA-ն
իրական ժամանակում գեներացնում ու մշակում է սենսորային տվյալները, իսկ ծրագրային
հարթակը դրանք պահպանում, վերլուծում է և հետադարձ կապով փոխում FPGA-ի պարամետրերը՝
առանց նոր bitstream գեներացնելու։}

Սկզբնական իրականացումը \textbf{ֆիզիկական սենսորներ չի պահանջում}. սենսորի վարքը
կարող է էմուլացվել անմիջապես FPGA-ի ներսում՝ կոնֆիգուրելի RTL ազդանշան-գեներատոր
մոդուլներով։ Ֆիզիկական սենսորները կարող են միանալ հետագայում՝ որպես ընդլայնում։

\section{Համակարգի ընդհանուր ճարտարապետություն}

\adddiagram[fig:arch]{architecture}{0.60}{Համակարգի ընդհանուր ճարտարապետությունը՝ AUP ZU3-ի PL/PS հատվածներից մինչև Web հարթակ.}

Ինչպես ցույց է տրված Նկ.~\ref{fig:arch}-ում, FPGA-ի ծրագրավորվող տրամաբանությունը
(PL) գեներացնում և մշակում է ազդանշանները, ապա \textbf{AXI}-ի միջոցով տվյալները
փոխանցվում են մշակող համակարգին (PS՝ ARM Cortex-A53)։ PS-ը \textbf{Ethernet}-ով
կապվում է սերվերի հետ, որտեղ գործում են Backend/API-ն, տվյալների բազան և Web
հարթակը (վիզուալիզացիա, վերլուծություն, մոդելավորում, կոնֆիգուրացիա)։
